\section{A resource-transition reading of the datum}\label{sec:benchmark}

The witness is abstract by design. This section gives it a Schaefer-style biomass reading using synthetic \emph{model biomass units} (MBU), not a calibrated fishery model. Biomass dynamics are specified, but the income-margin and reserve coordinates are normalized illustrative states whose paths are stipulated by the witness; they are not calculated from prices, costs, or fleet accounts. Northern-cod data enter only through a separately declared historical scale comparison (full detail in SI~S5). No model parameter is estimated from those data, and no fit, forecast, validation, or empirical performance claim is made.

\paragraph{State variables and scales} Let $B_m(t)$ be biomass in MBU over one annual review interval $t\in[0,1]$, with Schaefer surplus production $\varsigma(B_m)=rB_m(1-B_m/K)$, $r=4\,\mathrm{yr}^{-1}$, $K=10$ MBU, and model floor $B_{\mathrm{lim}}=2$ MBU. The first assessment coordinate is the normalized biomass margin $s_1=(B_m-B_{\mathrm{lim}})/(1\,\mathrm{MBU})$. The second is a separate normalized income-margin coordinate $s_2=\Delta\Pi/\Pi_0$, for a declared positive reference scale $\Pi_0$; no revenue or cost law is supplied, so its paths are illustrative rather than empirically generated. The reserve $x$ is likewise in normalized fund units, with staged-plan cost $c=1$ fund unit. The heatwave is a one-time biomass shock of $\delta_0=1/2$ MBU at $t=1/2$; it does not directly disturb $s_2$. The dashboard weights act on these dimensionless normalized margins, not on raw tonnes and currency.

\paragraph{The schedule and its status} At the witness point $(x,s_1,s_2)=(1/2,6/5,6/5)$, the opening biomass is $B_m(0)=16/5$ MBU. The \act{fast} schedule follows the benign biomass line $B_{\mathrm{line}}(t)=16/5-3t$ for $0\le t\le1/2$ under
\[
\dot B_m=\varsigma(B_m)-H_m(t),\qquad H_m(t)=3+\varsigma(B_{\mathrm{line}}(t)),
\]
then closes the fishery ($H_m=0$). The quota falls from $1463/125\approx11.704$ to $2161/250\approx8.644$ MBU/yr, below the stipulated capacity $H_{\max}=13$ MBU/yr. At the midpoint, the benign biomass is $17/10$ MBU; the heatwave branch is reduced by $1/2$ MBU to $6/5$ MBU. The corresponding certified floor-margin tubes are $s_1:6/5\to-3/10\to6/5$ without the strike and $s_1:6/5\to-4/5\to6/5$ with it. The \act{slow} plan applies the analogous prescribed excursion to the normalized income margin, while holding the biomass coordinate fixed. \act{staged} spends one fund unit and increases each normalized margin by $e=1/4$. These income and reserve paths are stipulated parts of the datum; they are not derived from an economic submodel.

The zero-harvest recovery is logistic, not the straight line used as the certified tube envelope. From a midpoint biomass $b$, the nonlinear endpoint is
\[
B_m(1)=\frac{10}{1+(10/b-1)e^{-2}},
\]
which is about $6.02136$ MBU from $b=1.7$ and $5.01892$ MBU from $b=1.2$; neither returns to the abstract successor $16/5$ MBU. The piecewise-linear recovery envelope is conservative: up to $B_m=5$, surplus is at least $\varsigma(6/5)=528/125>4$ MBU/yr, exceeding the envelope's recovery rate, and after the actual stock reaches $5$ it remains above the envelope, which never exceeds $16/5$. Thus the nonlinear trajectory remains above the stated lower envelope, but matching the witness's endpoint requires a separately specified post-interval control or reset. The exact-tube theorem applies to the abstract datum; this biomass calculation alone does not establish the full two-coordinate transition or its terminal condition.

\paragraph{The dashboard reading} On the heatwave branch of \act{fast}, the normalized margin sum takes the values $(12/5,2/5,12/5)$, so the equal-weight dashboard remains above its zero threshold even though the biomass margin is negative at mid-interval. The licensing thresholds at the witness point are $\rho_1=2/3$ and $\rho_2=3/2$: each weight has at least one licensed plan, and the ranges overlap; they do not assign a unique or necessarily different plan to every weight. A biomass-focused weight rejects the \act{fast} trajectory, but may license \act{slow}, which protects biomass in this datum. The conclusion is therefore state-level, weight-adaptive acceptance, not acceptance of one fixed trajectory under every weight. The mid-interval event is a biomass-limit breach in the model, not a claim of biological extinction or collapse. A reserve top-up is a separate remedy: for a typed-failure state with $x<1$, the witness threshold is $\kappa^*=1-x$.

\begin{table}[!htbp]
\centering\small
\caption{The witness coordinates read as quantities in the illustrative benchmark.}
\label{tab:fisheries}
\begin{tabularx}{\textwidth}{@{}>{\raggedright\arraybackslash}p{0.3\textwidth}X@{}}
\toprule
Object & Illustrative reading \\
\midrule
Typed margins $s_i\ge0$ & normalized biomass margin $s_1$ and separately normalized income-margin coordinate $s_2$ \\
Scalarized operator $E_w$ & the declared dashboard $w\cdot(s_1,s_2)$; weights do not multiply raw tonnes and currency \\
Licensing thresholds $\rho_1,\rho_2$ & weight-ratio ranges over which the pulse (\act{fast}) or gradual (\act{slow}) plan is licensed \\
Acceptance gap $\FP$ & dashboard-certified yet typed-inadmissible transitions under the stated plan and disturbance class \\
Rescue set $R$; threshold $\kappa^*$ & states served by the staged plan and the normalized fund increment needed to reach its cost \\
Exact-tube semantics & checking every state in the stated tube, including the midpoint disturbance \\
Critical elasticity $\sigma^*$ & threshold in the constructed normalized-margin CES family; not a fitted biological parameter \\
\bottomrule
\end{tabularx}
\end{table}

\paragraph{Historical scale comparison, not calibration} DFO Science Advisory Report 2016/026 reports a Northern-cod LRP of $884.6$ kt for the cited assessment vintage; it is the rounded 1983--1989 mean in that series (Table A2, PDF p.~16), not an asserted current-2026 threshold. If, solely for a dimensional comparison, $2$ MBU is mapped to $884.6$ kt, then $1$ MBU is $442.3$ kt. This would map the witness opening biomass $3.2$ MBU to $1415.36$ kt, the adverse trough $1.2$ MBU to $530.76$ kt, $K=10$ MBU to $4423$ kt, the $0.5$-MBU heatwave to $221.15$ kt, and the capacity $13$ MBU/yr to $5749.9$ kt/yr. These large implied capacities underscore that the mapping is a unit anchor, not a plausible or fitted Northern-cod production model. DFO's assessed 2015 level was about $298.65$ kt, or $33.8\%$ of that historical LRP. The CES input $\lambda_1=1+s_1$ is a constructed normalized-margin index, not $B_m/B_{\mathrm{lim}}$: under this mapping, $B_m/B_{\mathrm{lim}}=1+s_1/2$, whereas $\lambda_1=1+s_1$. The income coordinate has its own normalization, and neither it nor the reserve has a physical currency scale here. The comparison makes no claim that the benchmark parameters or management schedules describe a real fishery.

\begin{figure}[!htbp]
\centering
\includegraphics[width=0.96\linewidth]{fig_benchmark.png}
\caption{The witness datum read as an illustrative resource transition. (a) The constructed normalized index $s_1+s_2$ remains positive while the biomass margin breaches its floor. (b) The quota schedule is plotted in model biomass units per year (MBU/yr); the fund coordinate is normalized and carries no currency unit. The piecewise-linear tube and the nonlinear logistic recovery have different endpoint semantics, as detailed above and in SI~S5. All curves are schematic summaries of the declared datum, not empirical trajectories.}
\label{fig:benchmark}
\end{figure}
